
%%%%%%%%%%%%%%%%%%%%%%%%%%%%%%%%%%%%%%%%%
% Beamer Presentation
% LaTeX Template
% Version 1.0 (10/11/12)
%
% This template has been downloaded from:
% http://www.LaTeXTemplates.com
%
% License:
% CC BY-NC-SA 3.0 (http://creativecommons.org/licenses/by-nc-sa/3.0/)
%
%%%%%%%%%%%%%%%%%%%%%%%%%%%%%%%%%%%%%%%%%

%----------------------------------------------------------------------------------------
%	PACKAGES AND THEMES
%----------------------------------------------------------------------------------------

\documentclass{beamer}

\mode<presentation> {

% The Beamer class comes with a number of default slide themes
% which change the colors and layouts of slides. Below this is a list
% of all the themes, uncomment each in turn to see what they look like.

%\usetheme{default}
%\usetheme{AnnArbor}
%\usetheme{Antibes}
%\usetheme{Bergen}
%\usetheme{Berkeley}
%\usetheme{Berlin}
%\usetheme{Boadilla}
%\usetheme{CambridgeUS}
%\usetheme{Copenhagen}
%\usetheme{Darmstadt}
%\usetheme{Dresden}
%\usetheme{Frankfurt}
%\usetheme{Goettingen}
%\usetheme{Hannover}
%\usetheme{Ilmenau}
%\usetheme{JuanLesPins}
%\usetheme{Luebeck}
%\usetheme{Madrid}
%\usetheme{Malmoe}
%\usetheme{Marburg}
%\usetheme{Montpellier}
%\usetheme{PaloAlto}
\usetheme{Pittsburgh}
%\usetheme{Rochester}
%\usetheme{Singapore}
%\usetheme{Szeged}
%\usetheme{Warsaw}

% As well as themes, the Beamer class has a number of color themes
% for any slide theme. Uncomment each of these in turn to see how it
% changes the colors of your current slide theme.

%\usecolortheme{albatross}
\usecolortheme{beaver}
%\usecolortheme{beetle}
%\usecolortheme{crane}
%\usecolortheme{dolphin}
%\usecolortheme{dove}
%\usecolortheme{fly}
%\usecolortheme{lily}
%\usecolortheme{orchid}
%\usecolortheme{rose}
%\usecolortheme{seagull}
%\usecolortheme{seahorse}
%\usecolortheme{whale}
%\usecolortheme{wolverine}

%\setbeamertemplate{footline} % To remove the footer line in all slides uncomment this line
%\setbeamertemplate{footline}[page number] % To replace the footer line in all slides with a simple slide count uncomment this line

%\setbeamertemplate{navigation symbols}{} % To remove the navigation symbols from the bottom of all slides uncomment this line
}

\usepackage{graphicx} % Allows including images
\usepackage{booktabs} % Allows the use of \toprule, \midrule and \bottomrule in tables

\usepackage{fontspec}
\setsansfont{Calibri}[
    Path=Font/,
    Scale=0.9,
    Extension = .ttf,
    UprightFont=*-Regular,
    BoldFont=*-Bold,
    ItalicFont=*-Italic,
    BoldItalicFont=*-BoldItalic
    ]
\usepackage{unicode-math}
\setmathfont{Libertinus Math}% 使用默认的数学字体或你偏好的数学字体 Latin Modern Math/Libertinus Math

\usepackage{ragged2e}
\justifying\let\raggedright\justifying %两端对齐

 \usepackage{tcolorbox} 
%----------------------------------------------------------------------------------------
%	TITLE PAGE
%----------------------------------------------------------------------------------------

\title[Asset Pricing]{\textbf{Asset Pricing}} % The short title appears at the bottom of every slide, the full title is only on the title page

\subtitle{Habit Formation}

\author[Cheng Hang]{Cheng Hang} % Your name
\institute[DUFE] % Your institution as it will appear on the bottom of every slide, may be shorthand to save space
{School of Fintech \\
Dongbei University of Finance \& Economics \\ % Your institution for the title page
\medskip
\textbf{chenghangch@qq.com}% Your email address
}
\date{\today} % Date, can be changed to a custom date

	\AtBeginSection[]
{
	\begin{frame}{Content}
		\tableofcontents[sectionstyle=show/shaded,subsectionstyle=show/shaded/hide] %这几个参数我也不知道该如何准确地解释，反正它们最终的效果是突出显示当前章节，而其它章节都进行了淡化处理
		\addtocounter{framenumber}{-1}  %目录页不计算页码
	\end{frame}
}

%----------------------------------------------------------------------------------------
%	Begin
%----------------------------------------------------------------------------------------

\begin{document}

\begin{frame}
\titlepage % Print the title page as the first slide
\end{frame}

\begin{frame}
    \frametitle{5.2 Campbell-Cochrane (1999) Model: Introduction}
    \begin{itemize}
        \item \textbf{Primary Goal:} Explain the equity volatility puzzle (large equity return volatility with low consumption growth volatility).
        \item Uses an absolute deviation of consumption from the habit level.
        \item This allows generating time-varying risk aversion even with homoskedastic consumption growth: $\Delta c_{t+1} = g + \epsilon_{c,t+1}$.
        \item Preferences:
        \[ u(C_t) = \frac{(C_t - X_t)^{1-\gamma}}{1-\gamma} \]
        \item Define the \textbf{surplus ratio} $S_t = \frac{C_t - X_t}{C_t}$ (how high consumption is relative to habit).
        \item \textbf{Relative Risk Aversion:}
        \[ -\frac{C_t u''(C_t)}{u'(C_t)} = \gamma \frac{C_t}{C_t - X_t} = \gamma \frac{1}{S_t} \]
        \item Effective risk aversion depends on the surplus ratio.
    \end{itemize}
\end{frame}

\begin{frame}
    \frametitle{5.2 Campbell-Cochrane (1999) Model: Habit Adjustment and SDF}
    \begin{itemize}
        \item \textbf{Surplus Ratio Adjustment:} They work with the path for log surplus consumption $s_t = \log(S_t)$:
        \[ s_{t+1} = (1-\phi)\bar{s} + \phi s_t + \lambda(s_t)\epsilon_{c,t+1} \]
        \item The term $\lambda(s_t)$ specifies the pass-through of consumption to the habit. It ensures $X_t$ doesn't fall below $C_t$ (where utility is undefined).
        \item This term is also responsible for time-varying volatility of the SDF, which amplifies risk premium movements.
        \item This equation implies that current habit is a function of all past consumptions:
        \[ x_t = \alpha + (1-\phi) \sum_{j=0}^{\infty} \phi^j c_{t-j} \]
        so habit responds slowly and linearly to log consumption.
        \item \textbf{SDF for CC Model:}
        \begin{align*}
            u'(C_t) &= (C_t - X_t)^{-\gamma} = C_t^{-\gamma} \left( \frac{C_t - X_t}{C_t} \right)^{-\gamma} = C_t^{-\gamma} S_t^{-\gamma} \\
            M_{t+1} &= \delta \left( \frac{S_{t+1}}{S_t} \right)^{-\gamma} \left( \frac{C_{t+1}}{C_t} \right)^{-\gamma} \\
            \implies \tilde{m}_{t+1} &= -\gamma\tilde{s}_{t+1} - \gamma\tilde{c}_{t+1} = -\gamma\lambda(s_t)\epsilon_{c,t+1} - \gamma\epsilon_{c,t+1}
        \end{align*}
    \end{itemize}
\end{frame}

\begin{frame}
    \frametitle{5.2 Campbell-Cochrane (1999) Model: Risk-Free Rate}
    \begin{itemize}
        \item Risk-free rate calculated in the usual way: $r_{f,t+1} = -E_t[m_{t+1}] - \frac{1}{2}\text{var}_t(m_{t+1})$.
        \begin{align*}
            r_{f,t+1} &= -E_t[\log(\delta) - \gamma\Delta s_{t+1} - \gamma\Delta c_{t+1}] - \frac{1}{2}\text{var}_t(-\gamma\tilde{s}_{t+1} - \gamma\tilde{c}_{t+1}) \\
            &= -\log(\delta) + \gamma E_t[s_{t+1} - s_t] + \gamma g - \frac{\gamma^2\sigma_c^2}{2} (1+\lambda(s_t))^2 \\
            &= -\log(\delta) + \gamma g + \gamma(1-\phi)(\bar{s} - s_t) - \frac{\gamma^2\sigma_c^2}{2} (1+\lambda(s_t))^2
        \end{align*}
        \item Three main terms:
        \begin{enumerate}
            \item $-\log(\delta) + \gamma g$: Standard CCAPM terms.
            \item $\gamma(1-\phi)(\bar{s} - s_t)$: When $s_t$ is low (times are bad), consumption is closer to habit. Over time, habit adjusts, so the agent gets used to bad times. This makes them want to borrow against the future, driving up interest rates.
            \item $-\frac{\gamma^2\sigma_c^2}{2} (1+\lambda(s_t))^2$: In CC parametrization, $\lambda(s_t)$ declines in $s_t$. In a recession (low $s_t$), volatility is large, leading to elevated precautionary savings demand, driving interest rates down.
        \end{enumerate}
        \item In CC parametrization, the last two effects exactly offset each other, implying a constant riskless rate.
    \end{itemize}
\end{frame}

\begin{frame}
    \frametitle{5.2 Campbell-Cochrane (1999) Model: Other Results}
    \begin{itemize}
        \item \textbf{Sensitivity function $\lambda(s_t)$:} Parametrized such that $\lambda(s_t) \to \infty$ as $s_t \to 0$. This is needed so habit is very sensitive to consumption, preventing consumption from falling below habit. This implies that as $s_t \to 0$, the volatility of SDF $\to \infty$.
        \item \textbf{Habit Predetermination:} Habit is predetermined in the steady state ($dx/dc = 0$ at $s_t = \bar{s}$) and near the steady state ($d(dx/dc)/ds = 0$ at $s_t = \bar{s}$). $dx/dc$ is a U-shaped function of $s_t$, tangent to 0 at $s_t = \bar{s}$.
        \item \textbf{Calibration:} Model parameters (e.g., $\phi$) are calibrated to arrive at a steady state surplus $S = 0.057$, meaning habit is 94\% of consumption.
        \item \textbf{Effective Risk Aversion:} The steady state risk aversion is $\gamma/S = 2/0.057 \approx 30$. This means the model resembles a power utility model with high risk aversion, and thus doesn't solve the equity premium puzzle on its own.
        \item \textbf{Predictable Stock Prices:} Since risk aversion moves with consumption, it generates large predictable movements in stock prices.
        \item \textbf{Countercyclical Risk Premium:} After consumption is lowered closer to the habit level, effective risk aversion increases, driving the risk premium up.
    \end{itemize}
\end{frame}


\end{document}
